\uuid{vCzC}
\titre{Nombres premiers}
\auteur{bodin ; croizat ; sacre}
\organisation{exo7}
\licence{CC BY-NC-SA 4.0}
\datecreate{2022-03-08}
\niveau{L1}
\module{Algèbre}
\chapitre{Arithmétique dans Z}
\sousChapitre{Nombres premiers, nombres premiers entre eux}
\difficulte{2}
\theme{Nombres premiers}
\exo7qcmid{hilisit.arithmetique.034}

\contenu{
\begin{question}
\qclassification{103.04}
 Soit $p \geq 3$ un nombre premier et $ p = 4q + r$ le résultat de sa division euclidienne par $4$. On peut alors avoir :
    \begin{answers}
        \bad{$r=0$}
        \good{$r=1$}
        \bad{$r=2$}
        \good{$r=3$}
    \end{answers}
    \begin{explanations}
    Un nombre premier $\geq 3$ est nécessairement impair. Ceci exclut donc les possibilité $r=0$ et $r=2$ qui correspondent à des nombres pairs.\\
    On peut à titre d'exemple obtenir pour $p=3$ que $r=3$ ; et pour $p=5$ que  $r=1$.
    \end{explanations}
\end{question}
}
